\documentclass{article}
\usepackage{apstats-poster}
% Formula IDs: z-test-stat, p-value-interp, type-i-error, type-ii-error, power
\colorlet{PosterAccent}{UnitOrange}

\begin{document}
\begin{posterpage}
\PosterHeader[82pt]{UnitOrange}{3}{INFERENCE FOR PROPORTIONS \& CATEGORICAL DATA}
  {P10}{SIGNIFICANCE TESTS}
  {master recipe \quad\textbullet\quad also used in Units 4 and 5}
\InferenceSubstripes

\PosterZone{4.72in}{15.15in}{ZONE A · THE FORMULAS + THE FOUR-STEP FRAME}{%
  \FormulaCardSized{47pt}{z=\frac{\text{statistic}-\text{parameter under }H_0}{\text{standard error of the statistic}}}
  \vspace{.13in}
  {\fontsize{26pt}{33pt}\selectfont
    SE estimates the statistic's sampling SD; use the test's null model.
    A \(t\)-procedure uses \(t\) with its \textbf{df}.}
  \vspace{.19in}
  \FormulaCardSized{38pt}{P\text{-value}=P(\text{observed or more extreme statistic}\mid H_0\text{ true})}
  \vspace{.25in}
  \begin{minipage}[t]{12.0in}
    \MiniHeading{THE SAME FOUR STEPS AS A CONFIDENCE INTERVAL}\par\vspace{.13in}
    \StepCard{1}{STATE}{Define the parameter. Give \(H_0\), \(H_a\), and \(\alpha\) in context.}
    \StepCard{2}{PLAN}{Name the test. Check Random \(\to\) 10\% \(\to\) Normal / Large Counts.}
    \StepCard{3}{DO}{Statistic, df if needed, and P-value. Draw and shade the curve.}
    \StepCard{4}{CONCLUDE}{Compare P with \(\alpha\); state evidence for \(H_a\) in context.}
  \end{minipage}\hfill
  \begin{minipage}[t]{7.65in}
    \MiniHeading{THE DECISION}\par\vspace{.13in}
    \TextCard{\fontsize{30pt}{39pt}\selectfont
      \(\boldsymbol{P\le\alpha}\)\par
      \(\longrightarrow\) reject \(H_0\)\par
      \textbf{Convincing evidence} for \(H_a\).}
    \vspace{.20in}
    \TextCard{\fontsize{30pt}{39pt}\selectfont
      \(\boldsymbol{P>\alpha}\)\par
      \(\longrightarrow\) fail to reject \(H_0\)\par
      \textbf{Not convincing evidence} for \(H_a\).}
  \end{minipage}
}

\PosterZone{15.30in}{21.30in}{ZONE B · ERRORS + POWER}{%
  \begin{minipage}[t]{8.0in}
    \vspace{0pt}\RaggedRight
    \MiniHeading{DECISION × TRUE STATE}\par\vspace{.14in}
    {\fontsize{25pt}{32pt}\selectfont
      \renewcommand{\arraystretch}{1.50}
      \begin{tabularx}{\linewidth}{>{\RaggedRight\arraybackslash}p{2.45in}|>{\centering\arraybackslash}X|>{\centering\arraybackslash}X}
        & \(H_0\) true & \(H_a\) true\\\hline
        Reject \(H_0\) & Type I error & Correct\\\hline
        Fail to reject \(H_0\) & Correct & Type II error
      \end{tabularx}}
    \vspace{.18in}
    {\fontsize{26pt}{33pt}\selectfont
      \(\alpha\): significance level.\par
      \(\beta\) and power depend on a specific true effect.}
  \end{minipage}\hfill
  \begin{minipage}[t]{11.65in}
    \vspace{0pt}\RaggedRight
    \FormulaCardSized{29pt}{\begin{gathered}
      \alpha=P(\text{Type I error})\\[6pt]
      =P(\text{reject }H_0\mid H_0\text{ true})
    \end{gathered}}
    \vspace{.12in}
    \FormulaCardSized{29pt}{\begin{gathered}
      \beta=P(\text{Type II error})\\[6pt]
      =P(\text{fail to reject }H_0\mid H_a\text{ true})
    \end{gathered}}
    \vspace{.12in}
    {\fontsize{32pt}{40pt}\selectfont
      \(\textbf{Power}=1-\beta\)}\par
    {\fontsize{25pt}{32pt}\selectfont
      Increases with larger \(n\), larger \(\alpha\), or a larger true effect,
      with other factors held fixed.}
  \end{minipage}
}

\PosterZone{21.45in}{24.35in}{ZONE C · SAY IT IN WORDS}{%
  \SentenceFrame{“Because the P-value of \PosterBlank{1.0in} is [less/greater]
    than \(\alpha=\) \PosterBlank{1.0in}, we [reject / fail to reject] \(H_0\).
    There [is / is not] convincing evidence that \PosterBlank{4.2in}.”}
  \vspace{.10in}
  {\fontsize{24pt}{30pt}\selectfont
    Fill the last blank with \(H_a\) in context. At \(P=\alpha\), use “equal to” and reject.}
}

\PosterZone{24.50in}{27.45in}{ZONE D · WATCH OUT}{%
  \begin{minipage}[t]{6.10in}
    \MiniHeading{KEEP THE DECISION PRECISE}\par\vspace{.10in}
    {\fontsize{27pt}{35pt}\selectfont
      Say “fail to reject \(H_0\).”\par
      Never “accept \(H_0\).”}
  \end{minipage}\hfill
  \begin{minipage}[t]{6.10in}
    \MiniHeading{NAME THE PARAMETER}\par\vspace{.10in}
    {\fontsize{27pt}{35pt}\selectfont
      Hypotheses use \(p\) or \(\mu\),\par
      not \(\hat p\) or \(\bar x\).}
  \end{minipage}\hfill
  \begin{minipage}[t]{7.15in}
    \MiniHeading{CHOOSE TAILS BEFORE DATA}\par\vspace{.10in}
    {\fontsize{27pt}{35pt}\selectfont
      One-sided vs two-sided changes P. Let the question choose the tail.}
  \end{minipage}
}
\WorksheetTag{u6\_l4--7, l10--11 · u7\_l4--5, l8--10}
\end{posterpage}
\end{document}
